% SPDX-FileCopyrightText: Copyright (c) 2024-2026 Yegor Bugayenko
% SPDX-License-Identifier: MIT

\section{Math}\label{sec:math}

All objects in this Section belong to the \ff{ms} package.
%
The \deff{random} object is a pseudo-random number generator
  parameterized by a \deff{seed}; when dataized it behaves as a
  \deff{number} between zero and one.
The \deff{next} attribute is the next generator in the sequence,
  the \deff{clocked} attribute seeds the generator from a given clock,
  and the \deff{pseudo} attribute seeds it from the system clock.

The \deff{angle} object is a decorator of \deff{number}.
It assumes that the angle is in radians and has the following attributes
  for trigonometric functions:
  \adeff{sin}, \adeff{cos}, \deff{tan}, and \deff{ctan}.
The \deff{degrees} attribute converts the angle from radians to degrees,
  and the \deff{radians} attribute converts it back.

The constants \deff{pi} and \deff{e} approximate \(\pi\) and Euler's number.
%
The \deff{real} object is a decorator of \deff{number}.
%
The \deff{numbers} object decorates a \deff{tuple} of numbers and provides
  \deff{max} and \deff{min} attributes.
%
The \deff{integral} object calculates the definite integral of a function
  between two limits.
For example, this code approximates the integral of \(f(x) = x\)
  from one to ten with fifteen partitions, yielding approximately \(49.5\):

\begin{ffcode}
ms.integral
  x > [x]
  1
  10
  15
\end{ffcode}

A few more decorators of \deff{number}:

\begin{itemize}
    \item \deff{abs}: absolute value of the number
    \item \deff{exp}: Euler's number raised to the given power
    \item \deff{mod}: remainder of the division by another number
    \item \adeff{sqrt}: square root (\(\sqrt{\rho}\))
    \item \adeff{pow}: the number raised to a given power (\(\rho^x\))
    \item \adeff{ln}: natural logarithm (\(\ln \rho\))
    \item \adeff{acos}: arccosine
    \item \adeff{asin}: arcsine
\end{itemize}
